% !TeX program = lualatex
% !TeX encoding = utf8
% !TeX spellcheck = uk_UA
% !TeX root = ../main.tex

%% --------------------------------------------------------
%% Chapter file: Basics.tex
%% --------------------------------------------------------

\graphicspath{{\currfilebase/Pictures}}

\chapter{Плазма: від іонізованого газу до колективного середовища}\label{sec:plasma-basics}\hypertarget{plasma-basics}{}
\hypertarget{PlasmaPhysics}{}

\section{Що таке плазма?}\label{sec:what-is-plasma}
%% --------------------------------------------------------


\emphx{Плазма} (від грец. \emphx{plasma}~--- виліплене, оформлене)~--- це частково або
повністю іонізований газ, у~якому густини позитивних і~негативних зарядів практично
однакові, а~заряджені частинки поводяться \emphx{колективно}. Обидва виділені слова мають
точний зміст, який ми з'ясуємо поступово в~цьому розділі (квазінейтральність~--- у~§\,\ref{sec:quasineutrality},
колективність~--- у~§\,\ref{sec:debye-sphere}). Термін \enquote{плазма} запропонував
І.~Ленгмюр наприкінці 1920-х років; 1929~р. Л.~Тонкс і~Ленгмюр описали власні коливання
електронного газу, які ми розглянемо в~розд.~\ref{sec:plasma-waves}.

\emph{План розділу.} Спершу з'ясуємо, як виникає плазма (формула Саха) і~що означають
температури її компонент. Далі~--- головне: квазінейтральність, дебаївське екранування й~умови,
за яких плазма справді колективна (три умови: $N_D\gg1$, $\lambda_D\ll L$, $\omega\tau\gg1$; умова слабкого зв'язку $\Gamma\ll1$~--- це те саме, що $N_D\gg1$).
Наприкінці~--- застосування й~узагальнення: нелінійне екранування, шар біля стінки з~критерієм
Бома, зонд Ленгмюра та екранування у~виродженому газі.

Ступінь іонізації характеризують величиною
% ~~~~~~~~~~~~~~~~~~~~~~~~~~~~~~~~~~~~~~~~~~~~~~~~~~~~~~~~
\begin{equation}\label{eq:ionization-degree}
	\alpha = \frac{n_i}{n_i + n_n},
\end{equation}
% ~~~~~~~~~~~~~~~~~~~~~~~~~~~~~~~~~~~~~~~~~~~~~~~~~~~~~~~~
де $n_i$ і~$n_n$~--- концентрації іонів та нейтральних частинок. У~плазмі можуть одночасно
бути присутніми електрони, кілька сортів іонів і~нейтральні атоми чи молекули. Залежно від
$\alpha$ розрізняють \emphx{слабо іонізовану} (частки відсотка), \emphx{помірно
іонізовану} (кілька відсотків) і~\emphx{повністю іонізовану} ($\alpha\to1$) плазму.

\emphx{Чим плазма відрізняється від нейтрального газу?} Не просто наявністю заряджених
частинок. У~нейтральному газі частинки майже вільні й~обмінюються імпульсом
у~короткочасних зіткненнях. У~плазмі кулонівська сила спадає повільно ($\propto r^{-2}$): кількість частинок у~шарі
радіуса $r$ зростає як $r^2$, тож далекі частинки дають не менший внесок, ніж близькі,
і~кожна заряджена частинка \emph{одночасно} взаємодіє з~великою кількістю інших. Сила,
що діє на частинку, формується усередненим полем багатьох частинок, а~не парними
зіткненнями. Такі \enquote{колективні взаємодії} породжують екранування зарядів, плазмові
коливання, хвилі й~нестійкості. Тому плазма не зводиться до сукупності незалежних
частинок: її макроскопічні властивості визначаються спільною електромагнітною взаємодією.
Скільки саме частинок взаємодіє одночасно, визначає число $N_D$ у~сфері екранування
(§\,\ref{sec:debye-sphere}): кулонівська сила спадає повільно, але далі за дебаївський радіус поле заряду
екранується (§\,\ref{sec:debye-screening}), тож колективно діють лише частинки всередині цієї сфери.

Друга особливість~--- тісний зв'язок механіки з~електродинамікою. Плазма не лише рухається
в~зовнішньому електромагнітному полі, а~й~сама створює макроскопічні поля та струми, які
діють назад на її рух.

\begin{Theory}[Температура в~енергетичних одиницях]
Температуру $T$ вимірюємо в~енергетичних одиницях ($k_{\mathrm B}=1$, тож $k_{\mathrm B}T$ пишемо просто
як $T$); практично $T$ дають в~еВ, $1$~еВ $\approx11\,600$~К. Температура~--- це параметр
максвеллівського розподілу $f\propto\exp(-mv^2/2T)$; середня кінетична енергія поступального руху
$\langle mv^2/2\rangle=\tfrac32T$. Поки не сказано інакше, мова йде саме про такий розподіл.
Для оцінок зручні числа (СГС): $e=4{,}8\times10^{-10}$~СГС, $e^2\approx1{,}44\times10^{-7}$~еВ$\cdot$см,
$1$~СГС напруженості поля $=300$~В/см.
\end{Theory}

Плазма виникає різними способами: термічною іонізацією (нагрівання), електричним розрядом,
фотоіонізацією, ударною іонізацією, бомбардуванням частинками. Залежно від температури її
поділяють на \emphx{низькотемпературну} ($T\lesssim10^5$~К, тобто $\lesssim10$~еВ) і~\emphx{високотемпературну}
($T\gtrsim10^6$~К, тобто $\gtrsim100$~еВ; у~термоядерних установках~--- до $10^8$~К~$\approx10$~кеВ і~вище). Саме високотемпературна плазма потрібна для реакцій
\emphx{керованого термоядерного синтезу}. Діапазон реальних плазм надзвичайно широкий
(табл.~\ref{tab:plasma-types}).

\begin{table}[ht]
	\caption{Орієнтовні параметри різних плазм (значення наведено за порядком величини).}
	\label{tab:plasma-types}
	\centering\small
	\begin{platbl}{colspec={X[2.4,l,m]X[c,m]X[c,m]X[c,m]X[c,m]X[c,m]X[c,m]}}
		Плазма & $n_e$, $\mathrm{см}^{-3}$ & $T$, еВ & $\lambda_{De}$, см & $a=n^{-1/3}$, см & $N_D$ & $\Gamma$ \\
		Іоносфера & $10^{5}$ & $0{,}1$ & $0{,}74$ & $2{,}2\times10^{-2}$ & $1{,}7\times10^{5}$ & $7\times10^{-5}$ \\
		Міжзоряний газ & $10$ & $1$ & $2{,}4\times10^{2}$ & $0{,}46$ & $5\times10^{8}$ & $3\times10^{-7}$ \\
		Сонячна корона & $10^{8}$ & $10^{2}$ & $0{,}74$ & $2{,}2\times10^{-3}$ & $1{,}7\times10^{8}$ & $7\times10^{-7}$ \\
		Ядро Сонця & $10^{26}$ & $10^{3}$ & $2{,}4\times10^{-9}$ & $2{,}2\times10^{-9}$ & $5$ & $7\times10^{-2}$ \\
		Токамак & $10^{14}$ & $10^{4}$ & $7{,}4\times10^{-3}$ & $2{,}2\times10^{-5}$ & $1{,}7\times10^{8}$ & $7\times10^{-7}$ \\
		Сонячний вітер (1~а.о.) & $10$ & $10$ & $7{,}4\times10^{2}$ & $0{,}46$ & $1{,}7\times10^{10}$ & $3\times10^{-8}$ \\
		Газовий розряд & $10^{12}$ & $3$ & $1{,}3\times10^{-3}$ & $10^{-4}$ & $9\times10^{3}$ & $5\times10^{-4}$ \\
		Електрони в~металі & $10^{23}$ & вироджені, $E_F\sim5$ & --- & $2\times10^{-8}$ & --- & --- \\
	\end{platbl}
\end{table}

Концентрації в~таблиці відрізняються на 25 порядків (від $10$ до $10^{26}~\mathrm{см}^{-3}$). У~правих стовпцях наведено
похідні величини: дебаївську довжину $\lambda_{De}$, середню відстань між частинками $a=n^{-1/3}$, число частинок
у~дебаївській сфері $N_D$ і~параметр зв'язку $\Gamma$ (їх означено в~§\,\ref{sec:debye-sphere} і~§\,\ref{sec:coupling}). Для плазм із~$N_D\gg1$ маємо
$\lambda_{De}\gg a$, тобто колективність; ядро Сонця ($N_D\sim5$, $\Gamma\sim0{,}07$) лежить на межі застосовності (за $\Gamma$ плазма слабкозв'язана, але $N_D\sim5$~--- \enquote{багато} вже сумнівне, див.\ §\,\ref{sec:coupling}), а~електрони в~металі
\emphx{вироджені} (їхня енергія визначається не теплом, а~принципом Паулі: $T\ll E_F$) й~потребують квантового опису
(§\,\ref{sec:thomas-fermi}); $E_F$~--- енергія Фермі, тобто найбільша енергія електронів при $T=0$. Діаграму на рис.~\ref{fig:plasma-types} читають так: кожна точка~--- плазма з~даними $n$ і~$T$, а~лінії
відокремлюють області, де змінюється фізика ($\hbar$~--- стала Планка, $c$~--- швидкість світла).


\begin{figure}[ht]\centering
	\localinput{plasma_types.tikz}
	\caption{Діаграма $\lg n$--$\lg T$ видів плазми: лінії $N_D\sim1$ (орієнтовна межа слабкозв'язаної плазми; практично потрібно $N_D\gtrsim100$), $T=E_F$ (виродження)
	та $T=m_ec^2$ (релятивізм) обмежують область класичної ідеальної плазми $\max\{e^2n^{1/3},\hbar^2n^{2/3}/m_e\}\ll T\ll m_ec^2$}
	\label{fig:plasma-types}
\end{figure}


%% --------------------------------------------------------
\section{Іонізація: формула Саха}\label{sec:saha}
%% --------------------------------------------------------


Скільки енергії потрібно, щоб газ став плазмою? Енергія іонізації атома водню
$I=13{,}6$~еВ, що відповідає $1{,}6\times10^{5}$~К, проте газ іонізується помітно при
значно нижчих температурах. Пояснює це формула Саха, яка визначає ступінь іонізації газу
в~стані термодинамічної рівноваги.

\begin{Theory}[Нагадаємо зі статистичної фізики]
Для ідеального газу з~концентрацією $n$ хімічний потенціал
$\mu=T\ln\bigl(n\lambda_T^3/g\bigr)$, де
$\lambda_T=\sqrt{2\pi\hbar^2/mT}$~--- теплова довжина хвилі де~Бройля, $g$~--- статистична вага (кратність виродженого внутрішнього стану; для електрона зі~спіном $g_e=2$).
Умова хімічної рівноваги реакції $\text{атом}\rightleftharpoons\text{іон}+\text{електрон}$:
$\mu_a=\mu_i+\mu_e$. Енергія іонізації $I$ зсуває нуль енергії: атом лежить на $I$ нижче за
пару \enquote{іон + електрон}, тобто $\mu_a=T\ln(n_n\lambda_{Ta}^3/g_a)-I$ ($n_n$~--- концентрація атомів). Підставляючи, маємо
$\ln\dfrac{n_e\lambda_{Te}^3}{g_e}+\ln\dfrac{n_i\lambda_{Ti}^3}{g_i}=\ln\dfrac{n_n\lambda_{Ta}^3}{g_a}-\dfrac IT$,
звідки
$\dfrac{n_en_i}{n_n}=\dfrac{g_eg_i}{g_a}\,\dfrac{\lambda_{Ta}^3}{\lambda_{Te}^3\lambda_{Ti}^3}\,e^{-I/T}$.
Оскільки $m_a\approx m_i$, то $\lambda_{Ta}\approx\lambda_{Ti}$, і~залишається
$\dfrac{n_en_i}{n_n}=\dfrac{g_eg_i}{g_a}\,\lambda_{Te}^{-3}\,e^{-I/T}$.
Для водню $g_eg_i/g_a=2\cdot1/2=1$. Тут $T$~--- температура електронів, а~газ вважаємо ідеальним і~рівноважним.
(Врізку можна пропустити при першому читанні: далі важливий лише результат~--- формула Саха.)
\end{Theory}

Для водню формула Саха має вигляд
% ~~~~~~~~~~~~~~~~~~~~~~~~~~~~~~~~~~~~~~~~~~~~~~~~~~~~~~~~
\begin{equation}\label{eq:saha}
	\frac{n_e n_i}{n_n} = \left(\frac{m T}{2\pi\hbar^2}\right)^{3/2}\exp\left(-\frac{I}{T}\right)
	\approx 2{,}4\times10^{15}\,T^{3/2}\,e^{-I/T}~\mathrm{см}^{-3},
\end{equation}
% ~~~~~~~~~~~~~~~~~~~~~~~~~~~~~~~~~~~~~~~~~~~~~~~~~~~~~~~~
де в~числовій формі $T$ підставляють у~кельвінах, а~тоді й~$I$ в~показнику експоненти
треба виразити в~кельвінах: $I/k_{\mathrm B}\approx1{,}58\times10^{5}$~К (тут, і~лише тут, повертаємо $k_{\mathrm B}$ явно). Для інших
елементів у~правій частині з'являється відношення статистичних ваг $g_eg_i/g_a$.

Дві обставини визначають зміст формули. По-перше, експоненціальний множник: ступінь
іонізації стає помітним, коли $T$ лише у~$10$--$20$ разів менша за $I$.
По-друге, перед експонентою стоїть величина $n_q=(mT/2\pi\hbar^2)^{3/2}=\lambda_{Te}^{-3}$,
обернена до квантового об'єму електрона ($\lambda_{Te}$~--- теплова довжина хвилі де~Бройля
з~врізки вище). Об'єм $\lambda_{Te}^3$~--- це «клітинка», яку займає один квантовий стан електрона, тож $n_q=\lambda_{Te}^{-3}$~--- кількість
станів, доступних електрону в~одиниці об'єму (для $T=1$~еВ $n_q\approx3\times10^{21}~\mathrm{см}^{-3}$). Класична статистика
застосовна, поки $n_e\ll n_q$. Для типової розрідженої плазми
$n_q/n=1/(n\lambda_{Te}^3)\sim10^{6}$--$10^{10}$: на кожен електрон припадає величезна кількість вільних станів, тому
вільний стан значно вигідніший за зв'язаний в~атомі (виграш ентропії) і~перекриває
больцманівський фактор $e^{-I/T}$ уже при $T\sim I/10\ldots I/20$.
Звідси ще один наслідок: за однакової температури \emph{чим розрідженіший газ, тим легше він іонізується}~---
менша $n$ означає більше вільних станів на електрон.

Для контрасту візьмемо кімнатну температуру ($T\approx0{,}026$~еВ) і~газ з~$I=14{,}6$~еВ при $n\approx3\times10^{19}~\mathrm{см}^{-3}$
(атоми азоту; реальний $\mathrm N_2$ молекулярний, але порядок величини той самий): $e^{-I/T}\sim10^{-245}$, тож
ступінь іонізації $\alpha\sim10^{-122}$, тобто практично нуль. А~водень при $1$~еВ і~густині $10^{17}~\mathrm{см}^{-3}$ іонізований уже на~$\sim20\%$.
Щоб знайти $\alpha$, зауважимо: кожен акт іонізації народжує пару іон${}+{}$електрон, тож $n_e=n_i$.
Позначивши повну концентрацію частинок $n=n_i+n_n$, маємо $n_e=n_i=\alpha n$, $n_n=(1-\alpha)n$, і~рівняння~\eqref{eq:saha}
набуває вигляду $\alpha^2/(1-\alpha)=n_q\,e^{-I/T}/n$~--- квадратне рівняння відносно $\alpha$.
Для водню з~повною концентрацією частинок $n=10^{17}~\mathrm{см}^{-3}$
його розв'язок дає
% ~~~~~~~~~~~~~~~~~~~~~~~~~~~~~~~~~~~~~~~~~~~~~~~~~~~~~~~~
\begin{equation*}
	\alpha\approx1{,}3\times10^{-4}\ (0{,}5~\text{еВ}),\qquad
	\alpha\approx0{,}18\ (1~\text{еВ}),\qquad
	\alpha\approx0{,}99\ (2~\text{еВ}).
\end{equation*}
% ~~~~~~~~~~~~~~~~~~~~~~~~~~~~~~~~~~~~~~~~~~~~~~~~~~~~~~~~
Газ іонізується у~вузькому інтервалі температур (рис.~\ref{tikz:saha-curve}): для водню
з~повною концентрацією $10^{17}~\mathrm{см}^{-3}$~--- між $0{,}5$ і~$2$~еВ (при іншій
густині інтервал зсувається логарифмічно). Перехід \emphx{плавний}: $\alpha(T)$~--- гладка
функція, розриву немає, тобто це \emphx{кросовер}\footnote{Кросовер~--- плавний перехід між двома станами системи при зміні параметра
(тут температури): без стрибка властивостей і~без певної температури переходу, на відміну від фазового переходу,
як-от кипіння води}, а~не фазовий перехід. Так є в~розрідженому газі, де застосовна формула~\eqref{eq:saha}. У~дуже густій
речовині (сильнозв'язана плазма) можливий і~справжній фазовий перехід, але це вже окреме
питання. У~розрядах і~під дією випромінювання ступінь іонізації може значно відрізнятися
від рівноважного, тому формула~\eqref{eq:saha} дає лише орієнтир.

%---------------------------------------------------------
\begin{figure}[ht]\centering
\localinput{saha-curve.tikz}
\caption{Залежність ступеня іонізації водню від температури (формула Саха); пунктиром позначено $T=I/10\approx1{,}4$~еВ,
де іонізація вже помітна, хоча $T\ll I$}
\label{tikz:saha-curve}
\end{figure}


%% --------------------------------------------------------
\section{Температура плазми та два компоненти}\label{sec:two-temperatures}
%% --------------------------------------------------------


У~плазмі електрони й~іони можуть мати різні середні кінетичні енергії, тому взагалі треба
вводити дві температури $T_e$ та~$T_i$. У~фізиці плазми температуру часто виражають
в~енергетичних одиницях: \enquote{температура в~еВ} означає величину $T$
в~електронвольтах (співвідношення з~кельвінами $1$~еВ~$\approx11\,600$~К~--- у~врізці в~§\,\ref{sec:what-is-plasma}).

У~газорозрядній плазмі енергія від зовнішнього поля спершу передається електронам: струм, а~отже й~джоулеве тепло $\vect j\cdot\vect E$, переносять переважно вони, бо значно рухливіші за іони. Через
велику різницю мас ($m_i\geq1836\,m_e$) електрон при пружному зіткненні з~іоном передає
тому лише частку $\delta E/E\le4m_em_i/(m_e+m_i)^2\approx4m_e/m_i\approx0{,}2\%$ (для протона) своєї енергії~--- як тенісний м'ячик, що відскакує від вантажівки
(центральний пружний удар з~теоретичної механіки), тому часто $T_e>T_i$. Лише за достатньо
великої густини та тривалого часу електрон-іонні зіткнення вирівнюють температури:
$T_e=T_i=T$. Отже, вираз \enquote{температура плазми} без уточнення неоднозначний,
і~в~нерівноважній плазмі $T_e$ та~$T_i$ слід вказувати окремо. Якщо $T_e\neq T_i$, плазму
називають \emphx{двотемпературною}. Яскравий приклад~--- позитивний стовп
тліючого розряду (довга світна однорідна ділянка розряду в~трубці між катодом і~анодом)
при тиску $\sim0{,}1$~мм рт.~ст.: температура електронів досягає
$\sim5\times10^4$~К ($\approx4$~еВ), тоді як іони мають лише кілька сотень кельвінів.
Різні часи релаксації (розд.~\ref{sec:plasma-collisions}) пояснюють, чому така система існує
тривалий час: спочатку встановлюється максвеллівський розподіл електронів, потім іонів,
і~лише згодом вирівнюються їхні температури.


%% --------------------------------------------------------
\section{Квазінейтральність і~масштаб розділення зарядів}\label{sec:quasineutrality}
%% --------------------------------------------------------


Якби в~плазмі була присутня велика кількість нескомпенсованих зарядів, виникали б сильні
поля, які приводили б заряди в~рух. Тому на макроскопічних масштабах густини позитивного
й~негативного зарядів практично збігаються.

Нехай у~плазмі з~концентраціями $n_e$ електронів (заряд $-e$) та $n_i$ іонів (заряд $+Ze$)
виникла просторова локальна різниця густин. За теоремою Гаусса виникне кулонівське поле:
% ~~~~~~~~~~~~~~~~~~~~~~~~~~~~~~~~~~~~~~~~~~~~~~~~~~~~~~~~
\begin{equation}
	\nabla \cdot \vect{E} = 4\pi \rho = 4\pi e (Z n_i - n_e).
\end{equation}
% ~~~~~~~~~~~~~~~~~~~~~~~~~~~~~~~~~~~~~~~~~~~~~~~~~~~~~~~~
Оскільки питомий заряд електрона $e/m_e$ величезний, а~іони важкі й~повільні, це поле за дуже короткий час
створить потужні сили, які змістять електрони так, щоб знівелювати розділення зарядів. Характерний час~---
$\lambda_{De}/v_{Te}$, де $v_{Te}\sim\sqrt{T_e/m_e}$~--- теплова швидкість електронів; для $n=10^{12}~\mathrm{см}^{-3}$,
$T=1$~еВ це $\sim2\times10^{-11}$~с. Цей час є оберненою \emph{плазмовою частотою} $\omega_p^{-1}$ (її виведемо
в~розд.~\ref{sec:plasma-waves}); поки що важливо лише, що вирівнювання дуже швидке.
Оцінимо силу цього поля: плазма з~$n=10^{12}~\mathrm{см}^{-3}$ і~розміром $L=1$~см, розділення зарядів лише на 1\%, $\delta n=10^{10}~\mathrm{см}^{-3}$.
Тоді $E\sim4\pi e\,\delta n\,L\approx60$~СГС$\approx2\times10^{4}$~В/см, тоді як тепловому руху при $T=1$~еВ відповідає поле $T/(eL)\sim1$~В/см (поле суттєве, якщо його робота
$eEL$ на відстані $L$ порівнянна з~тепловою енергією $T$).
Навіть відхилення $\delta n/n\sim10^{-6}$ дає $E\sim2$~В/см, тобто порівнянне з~тепловим: розділення зарядів неможливе вже на рівні $10^{-6}$. Загальну оцінку дано наприкінці §\,\ref{sec:debye-screening}.

Тому на макроскопічних масштабах плазма є \emphx{квазінейтральною}:
\begin{equation}\label{eq:quasineutrality}
	\rho = \sum_s q_s n_s = e(Z n_i - n_e) \approx 0 \quad \Rightarrow \quad n_e \approx Z n_i.
\end{equation}
% ========================================================
\begin{Attention}
	Квазінейтральність~--- це не тотожний нуль у~кожній точці, а~макроскопічне наближення.
	Локальні відхилення від нейтральності можливі й~саме вони лежать в~основі екранування
	та плазмових коливань.
\end{Attention}
% ========================================================


%% --------------------------------------------------------
\section{Дебаївське екранування}\label{sec:debye-screening}
%% --------------------------------------------------------


Внесемо в~плазму малий пробний заряд $q$. Тут виникає перше принципове питання: якщо
кулонівська взаємодія далекодійна, чому поле окремого заряду не відчувається на довільно
великих відстанях? Відповідь полягає в~перерозподілі самих частинок плазми. Заряд
виявиться оточеним \enquote{шубою} частинок протилежного знака, і~його поле зменшиться
(рис.~\ref{tikz:debye-sphere}). Тепловий рух не дає такій \enquote{шубі} стиснутися до точки: вона розмазана на відстань порядку $\lambda_D$ (повний її заряд при цьому виявиться рівним $-q$, див.\ нижче). Знайдемо
характерну відстань, на якій відбувається це екранування.

%---------------------------------------------------------
\begin{SCfigure}[2][ht]\centering
\localinput{debye_sphere.tikz}
\caption{Пробний додатний заряд оточений електронною \enquote{шубою}, яка екранує його
поле}
\label{tikz:debye-sphere}
\end{SCfigure}
%---------------------------------------------------------

Кожен сорт частинок $s$ (заряд $q_s$, концентрація $n_s$) перебуває в~стаціонарному потенціалі $\varphi(\vect{r})$
у~тепловій рівновазі зі своєю температурою $T_s$, тому (з~розподілу Гіббса: ймовірність стану з~енергією $U$ пропорційна $e^{-U/T}$, а~потенційна енергія частинки в~полі $U=q_s\varphi$) розподіл частинок за Больцманом
% ~~~~~~~~~~~~~~~~~~~~~~~~~~~~~~~~~~~~~~~~~~~~~~~~~~~~~~~~
\begin{equation*}
	n_s=n_{s0}\exp\left(-\frac{q_s\varphi}{T_s}\right).
\end{equation*}
% ~~~~~~~~~~~~~~~~~~~~~~~~~~~~~~~~~~~~~~~~~~~~~~~~~~~~~~~~
Тут $n_{s0}$~--- незбурена концентрація вдалині від заряду, де $\varphi\to0$, а~сумарний
заряд плазми нульовий: $\sum_sq_sn_{s0}=0$. Густина заряду плазми
% ~~~~~~~~~~~~~~~~~~~~~~~~~~~~~~~~~~~~~~~~~~~~~~~~~~~~~~~~
\begin{equation*}
	\rho=\sum_sq_sn_{s0}\exp\left(-\frac{q_s\varphi}{T_s}\right).
\end{equation*}
% ~~~~~~~~~~~~~~~~~~~~~~~~~~~~~~~~~~~~~~~~~~~~~~~~~~~~~~~~
Для слабкого поля $|q_s\varphi|\ll T_s$ розкладемо експоненти. Нульовий
доданок зникає завдяки нейтральності, залишається
% ~~~~~~~~~~~~~~~~~~~~~~~~~~~~~~~~~~~~~~~~~~~~~~~~~~~~~~~~
\begin{equation*}
	\rho\simeq-\varphi\sum_s\frac{n_{s0}q_s^2}{T_s}.
\end{equation*}
% ~~~~~~~~~~~~~~~~~~~~~~~~~~~~~~~~~~~~~~~~~~~~~~~~~~~~~~~~
Підстановка в~рівняння Пуассона $\Delta\varphi=-4\pi(\rho+q\delta(\vect{r}))$ дає
(дельта-функція $\delta(\vect r)$ описує точковий пробний заряд і~забезпечує кулонівську поведінку потенціалу при $r\to0$)
% ~~~~~~~~~~~~~~~~~~~~~~~~~~~~~~~~~~~~~~~~~~~~~~~~~~~~~~~~
\begin{equation}\label{eq:poisson-screened}
	\Delta\varphi-\frac{\varphi}{\lambda_D^2}=-4\pi q\,\delta(\vect{r}),
	\qquad
	\tcbhighmath{\frac{1}{\lambda_D^2}=4\pi\sum_s\frac{n_{s0}q_s^2}{T_s}}.
\end{equation}
% ~~~~~~~~~~~~~~~~~~~~~~~~~~~~~~~~~~~~~~~~~~~~~~~~~~~~~~~~
Величина $\lambda_D$ називається \emphx{дебаївською довжиною} (радіусом екранування).

Поле пробного заряду сферично симетричне, $\varphi=\varphi(r)$, і~при $r>0$ маємо
$\Delta\varphi=\dfrac1r\dfrac{d^2(r\varphi)}{dr^2}$. Для $u=r\varphi$
рівняння~\eqref{eq:poisson-screened} спрощується до $u''=u/\lambda_D^2$, загальний
розв'язок якого $u=Ae^{-r/\lambda_D}+Be^{r/\lambda_D}$. Умова $\varphi\to0$ при
$r\to\infty$ виключає зростаючий доданок ($B=0$), а~близько до заряду потенціал має
збігатися з~кулонівським $q/r$, звідки $A=q$:
% ~~~~~~~~~~~~~~~~~~~~~~~~~~~~~~~~~~~~~~~~~~~~~~~~~~~~~~~~
\begin{equation}\label{eq:screened-potential}
	\varphi(r)=\frac qr\exp\left(-\frac r{\lambda_D}\right).
\end{equation}
% ~~~~~~~~~~~~~~~~~~~~~~~~~~~~~~~~~~~~~~~~~~~~~~~~~~~~~~~~
Кулонівське поле заряду експоненціально послаблюється на відстанях порядку $\lambda_D$
(рис.~\ref{pic:debye-screening}). Цей результат у~1923~р. одержали П.~Дебай і~Е.~Хюккель
для електролітів; плазма~--- його найчистіший прояв.

\begin{figure}[ht]\centering
	\localinput{potential.tikz}
	\caption{Дебаївське екранування пробного заряду: потенціал швидко спадає на відстанях
	порядку $\lambda_D$}
	\label{pic:debye-screening}
\end{figure}

Знайдемо заряд \enquote{шуби}. Густина заряду плазми
$\rho=-\Delta\varphi/4\pi=-\varphi/(4\pi\lambda_D^2)$ при $r>0$, тому
% ~~~~~~~~~~~~~~~~~~~~~~~~~~~~~~~~~~~~~~~~~~~~~~~~~~~~~~~~
\begin{equation*}
	\int\rho\,4\pi r^2dr=-\frac q{\lambda_D^2}\int_0^\infty re^{-r/\lambda_D}dr=-q.
\end{equation*}
% ~~~~~~~~~~~~~~~~~~~~~~~~~~~~~~~~~~~~~~~~~~~~~~~~~~~~~~~~
Отже, повний заряд оболонки в~точності компенсує $q$, а~її характерний радіус~---
$\lambda_D$. Для спостерігача на відстані $\gg\lambda_D$ пробний заряд разом з~оболонкою
практично нейтральний. Те саме стосується кожної зарядженої частинки самої плазми: для решти частинок вона \enquote{пробна}, тож усюди в~плазмі частинка оточена шубою завтовшки $\lambda_D$ (це використаємо в~§\,\ref{sec:debye-sphere}).


\subsection*{Скільки сортів частинок бере участь в~екрануванні?}


Для електронів і~одноразово заряджених іонів
% ~~~~~~~~~~~~~~~~~~~~~~~~~~~~~~~~~~~~~~~~~~~~~~~~~~~~~~~~
\begin{equation*}
	\frac1{\lambda_D^2}=\frac1{\lambda_{De}^2}+\frac1{\lambda_{Di}^2},
	\qquad
	\lambda_{De}=\sqrt{\frac{T_e}{4\pi n_ee^2}},
	\quad
	\lambda_{Di}=\sqrt{\frac{T_i}{4\pi n_ie^2}}.
\end{equation*}
% ~~~~~~~~~~~~~~~~~~~~~~~~~~~~~~~~~~~~~~~~~~~~~~~~~~~~~~~~~~~~~~
При $T_e=T_i=T$ маємо
% ~~~~~~~~~~~~~~~~~~~~~~~~~~~~~~~~~~~~~~~~~~~~~~~~~~~~~~~~
\begin{equation}\label{eq:debye-radius}
	\lambda_D=\sqrt{\frac{T}{8\pi ne^2}}=\frac{\lambda_{De}}{\sqrt2}.
\end{equation}
% ~~~~~~~~~~~~~~~~~~~~~~~~~~~~~~~~~~~~~~~~~~~~~~~~~~~~~~~~
Практично $\lambda_{De}\approx743\sqrt{T_e[\text{еВ}]/n_e[\mathrm{см}^{-3}]}$~см.

Формули вище припускають, що \emph{обидва} сорти встигають перебудуватися в~полі заряду, тобто що процес
повільний порівняно з~часом руху іонів на відстань $\lambda_D$ ($\sim\lambda_D/v_{Ti}$). Якщо процес швидкий,
іони лишаються практично на місцях, екранують лише електрони, і~замість $\lambda_D$ слід брати $\lambda_{De}$.
Зауважимо: у~двотемпературній плазмі з~$T_i\ll T_e$ сума дає $\lambda_D\approx\lambda_{Di}\ll\lambda_{De}$~---
холодні іони екранують \emph{статичний} заряд ефективніше за гарячі електрони, але лише якщо встигли досягти
рівноваги з~температурою $T_i$. У~літературі під \enquote{дебаївською довжиною} найчастіше розуміють саме
$\lambda_{De}$, тому в~кожній задачі слід перевіряти, які компоненти встигають перебудуватися.


\subsection*{Умова застосовності}


Лінеаризація експонент потребує $|q_s\varphi|\ll T_s$ принаймні на відстанях
$r\gtrsim\lambda_D$. Для $q=e$ на $r=\lambda_D$ це $e^2/\lambda_D\ll T$, або,
підставивши~\eqref{eq:poisson-screened}, $1/(8\pi n\lambda_D^3)\ll1$ (для двох сортів $T=8\pi ne^2\lambda_D^2$). Тобто теорія
виконується лише тоді, коли \emph{в~сфері радіуса $\lambda_D$ міститься багато частинок} (їхню кількість $N_D$ означено в~§\,\ref{sec:debye-sphere}; $1/(8\pi n\lambda_D^3)$~--- це, з~точністю до числового множника, $1/N_D$).
Це не технічна деталь, а~фізичний зміст колективного опису: екранування повинно
створюватися статистично усередненим полем багатьох частинок.

Зауважимо, що зовсім близько до пробного заряду ($r\lesssim e^2/T$) потенціал кулонівський і~$|q\varphi|\gtrsim T$,
тож лінеаризація там порушується. Але ця область набагато менша за $\lambda_D$ (їхнє відношення $e^2/(T\lambda_D)\sim1/(n\lambda_D^3)$~--- той самий малий
параметр $1/N_D$), тому на масштабі екранування формула~\eqref{eq:screened-potential} справедлива.

Тепер можна уточнити зміст квазінейтральності. Нехай на масштабі $L$ виникла локальна
неоднорідність концентрації $\delta n$. З~рівняння Пуассона $\varphi/L^2\sim4\pi e\,\delta
n$, а~потенціал не може суттєво перевищити $T/e$: за $e\varphi\gg T$ електрони з~тепловими енергіями $\sim T$ були б виштовхані або втягнені полем, і~відхилення зникло б. Тож $\varphi\lesssim T/e$ і
% ~~~~~~~~~~~~~~~~~~~~~~~~~~~~~~~~~~~~~~~~~~~~~~~~~~~~~~~~
\begin{equation*}
	\frac{\delta n}n\lesssim\frac{T}{4\pi ne^2L^2}=\left(\frac{\lambda_{De}}L\right)^2.
\end{equation*}
% ~~~~~~~~~~~~~~~~~~~~~~~~~~~~~~~~~~~~~~~~~~~~~~~~~~~~~~~~
Отже, квазінейтральність виконується тим точніше, чим менша дебаївська довжина порівняно
з~розміром плазми $L$:
% ~~~~~~~~~~~~~~~~~~~~~~~~~~~~~~~~~~~~~~~~~~~~~~~~~~~~~~~~
\begin{equation}\label{eq:size-condition}
	\lambda_D\ll L.
\end{equation}
% ~~~~~~~~~~~~~~~~~~~~~~~~~~~~~~~~~~~~~~~~~~~~~~~~~~~~~~~~
(У~нерівностях такого типу різницею між $\lambda_D$ і~$\lambda_{De}$ нехтуємо~--- вона лише в~множнику $\sqrt2$.) Біля стінок, електродів та інших меж ця умова порушується~--- там існують \emph{дебаївські
шари}, де квазінейтральності немає.


%% --------------------------------------------------------
\section{Дебаївська сфера. Умова колективності}\label{sec:debye-sphere}
%% --------------------------------------------------------


Електричне поле зарядженої частинки в~плазмі екранується іншими частинками. Щоб зрозуміти,
наскільки цей процес є колективним, розглянемо область простору радіуса $\lambda_{De}$
навколо даної частинки. Її називають \emphx{дебаївською сферою}.

Кількість електронів у~такій сфері дорівнює

% ~~~~~~~~~~~~~~~~~~~~~~~~~~~~~~~~~~~~~~~~~~~~~~~~~~~~~~~~
\begin{equation}\label{eq:ND}
    N_D=\frac{4\pi}{3}n_e\lambda_{De}^3.
\end{equation}
% ~~~~~~~~~~~~~~~~~~~~~~~~~~~~~~~~~~~~~~~~~~~~~~~~~~~~~~~~

Отже, $N_D$~--- це просто число частинок у~області, на масштабі якої відбувається
дебаївське екранування. Тут береться саме електронна довжина $\lambda_{De}$; заміна на
$\lambda_D=\lambda_{De}/\sqrt2$ змінила б $N_D$ лише в~$2^{3/2}\approx2{,}8$ раза, що для умови
$N_D\gg1$ несуттєво.

Якщо
% ~~~~~~~~~~~~~~~~~~~~~~~~~~~~~~~~~~~~~~~~~~~~~~~~~~~~~~~~
\begin{equation}\label{eq:ND-coll}
    N_D\gg1,
\end{equation}
% ~~~~~~~~~~~~~~~~~~~~~~~~~~~~~~~~~~~~~~~~~~~~~~~~~~~~~~~~
то в~екрануванні бере участь дуже багато частинок. У~цьому випадку не можна говорити про
взаємодію даної частинки з~одним конкретним сусідом: поле, що виникає в~результаті
перерозподілу зарядів, визначається сумарною дією великої кількості частинок. Екранування
має \emphx{колективний} характер.

Саме тому умова~\eqref{eq:ND-coll} є основною умовою колективності плазми.

Є й~спостережуваний наслідок. Парні зіткнення змінюють імпульс частинки з~частотою $\nu\sim\omega_p\ln N_D/N_D$
(орієнтовно; виведемо в~розд.~\ref{sec:plasma-collisions}; логарифм дають далекі проліти, а~не лише близькі зіткнення),
де $\omega_p=v_{Te}/\lambda_{De}$~--- плазмова частота, тобто швидкість
колективної відповіді (наступний розд.~\ref{sec:plasma-waves}; числово $\omega_p\approx5{,}6\times10^{4}\sqrt{n[\mathrm{см}^{-3}]}$~рад/с). Тому за $N_D\gg1$ колективні ефекти домінують над парними
зіткненнями: $\nu\ll\omega_p$. Саме це робить \enquote{усереднене поле багатьох частинок} головною силою в~плазмі.

Зрозуміємо тепер, коли ця умова виконується. Концентрація $n$ визначає характерну відстань
між частинками. Оскільки на одну частинку в~середньому припадає об'єм порядку $1/n$, маємо

% ~~~~~~~~~~~~~~~~~~~~~~~~~~~~~~~~~~~~~~~~~~~~~~~~~~~~~~~~
\begin{equation*}
    a^3\sim\frac{1}{n},
    \qquad
    a\sim n^{-1/3},
\end{equation*}
% ~~~~~~~~~~~~~~~~~~~~~~~~~~~~~~~~~~~~~~~~~~~~~~~~~~~~~~~~
де $a$~--- характерний масштаб міжчастинкової відстані.

Тоді

% ~~~~~~~~~~~~~~~~~~~~~~~~~~~~~~~~~~~~~~~~~~~~~~~~~~~~~~~~
\begin{equation*}
    N_D\sim\left(\frac{\lambda_D}{a}\right)^3.
\end{equation*}
% ~~~~~~~~~~~~~~~~~~~~~~~~~~~~~~~~~~~~~~~~~~~~~~~~~~~~~~~~

Отже, умова $N_D\gg1$ просто означає

% ~~~~~~~~~~~~~~~~~~~~~~~~~~~~~~~~~~~~~~~~~~~~~~~~~~~~~~~~
\begin{equation*}
    \lambda_D\gg a.
\end{equation*}
% ~~~~~~~~~~~~~~~~~~~~~~~~~~~~~~~~~~~~~~~~~~~~~~~~~~~~~~~~

Інакше кажучи, дебаївська довжина повинна бути набагато більшою за відстань між сусідніми
частинками. Тільки тоді в~області екранування міститься багато частинок і~екранування може
бути колективним.

\begin{Attention}
    Слабка взаємодія окремих частинок не означає відсутності колективних ефектів. Навпаки,
    колективність виникає саме тому, що слабкі внески великої кількості заряджених
    частинок складаються.
\end{Attention}


%% --------------------------------------------------------
\section{Слабкозв'язана та сильнозв'язана плазма}\label{sec:coupling}
%% --------------------------------------------------------


Є ще один спосіб оцінити, наскільки важливою є кулонівська взаємодія між сусідніми
частинками. Порівняємо потенціальну енергію їхньої взаємодії

% ~~~~~~~~~~~~~~~~~~~~~~~~~~~~~~~~~~~~~~~~~~~~~~~~~~~~~~~~
\begin{equation*}
    U_{\mathrm C}\sim\frac{e^2}{a}
\end{equation*}
% ~~~~~~~~~~~~~~~~~~~~~~~~~~~~~~~~~~~~~~~~~~~~~~~~~~~~~~~~

з~характерною енергією теплового руху

% ~~~~~~~~~~~~~~~~~~~~~~~~~~~~~~~~~~~~~~~~~~~~~~~~~~~~~~~~
\begin{equation*}
    U_T\sim T.
\end{equation*}
% ~~~~~~~~~~~~~~~~~~~~~~~~~~~~~~~~~~~~~~~~~~~~~~~~~~~~~~~~

Їхнє відношення називають \emphx{параметром зв'язку}:

% ~~~~~~~~~~~~~~~~~~~~~~~~~~~~~~~~~~~~~~~~~~~~~~~~~~~~~~~~
\begin{equation}\label{eq:Gamma}
    \Gamma=
    \frac{U_{\mathrm C}}{U_T}
    =
    \frac{e^2}{a\,T},
    \qquad
    a=n^{-1/3}.
\end{equation}
% ~~~~~~~~~~~~~~~~~~~~~~~~~~~~~~~~~~~~~~~~~~~~~~~~~~~~~~~~

Таким чином, $\Gamma$ безпосередньо показує, наскільки сильна кулонівська взаємодія
порівняно з~тепловим рухом.

При
% ~~~~~~~~~~~~~~~~~~~~~~~~~~~~~~~~~~~~~~~~~~~~~~~~~~~~~~~~
\begin{equation*}
    \Gamma\ll1
\end{equation*}
% ~~~~~~~~~~~~~~~~~~~~~~~~~~~~~~~~~~~~~~~~~~~~~~~~~~~~~~~~
теплова енергія значно більша за кулонівську. Частинки взаємодіють слабко, а~плазму
називають \emphx{слабкозв'язаною}. У~цьому режимі

% ~~~~~~~~~~~~~~~~~~~~~~~~~~~~~~~~~~~~~~~~~~~~~~~~~~~~~~~~
\begin{equation*}
    N_D\gg1,
\end{equation*}
% ~~~~~~~~~~~~~~~~~~~~~~~~~~~~~~~~~~~~~~~~~~~~~~~~~~~~~~~~
і~екранування має колективний характер.

Навпаки, при

% ~~~~~~~~~~~~~~~~~~~~~~~~~~~~~~~~~~~~~~~~~~~~~~~~~~~~~~~~
\begin{equation*}
    \Gamma\gtrsim1
\end{equation*}
% ~~~~~~~~~~~~~~~~~~~~~~~~~~~~~~~~~~~~~~~~~~~~~~~~~~~~~~~~
кулонівська взаємодія між сусідніми частинками стає порівнянною з~тепловою або перевищує
її. Частинки рухаються вже не як майже незалежні частинки, а~їхній рух суттєво
корельований. Таку систему називають \emphx{сильнозв'язаною}.

Для нашого визначення $a=n^{-1/3}$ зв'язок між $\Gamma$ та $N_D$ отримується безпосередньо
з~формули для дебаївської довжини. Справді, $T=e^2n^{1/3}/\Gamma$, тож
$\lambda_{De}^2=T/(4\pi ne^2)=1/(4\pi n^{2/3}\Gamma)$ і
$N_D=\tfrac{4\pi}{3}n\lambda_{De}^3=\tfrac{4\pi}{3}(4\pi)^{-3/2}\Gamma^{-3/2}$, тобто

% ~~~~~~~~~~~~~~~~~~~~~~~~~~~~~~~~~~~~~~~~~~~~~~~~~~~~~~~~
\begin{equation}\label{eq:ND-Gamma}
    N_D=
    \frac{1}{6\sqrt{\pi}}\,
    \Gamma^{-3/2}.
\end{equation}
% ~~~~~~~~~~~~~~~~~~~~~~~~~~~~~~~~~~~~~~~~~~~~~~~~~~~~~~~~

Тому
% ~~~~~~~~~~~~~~~~~~~~~~~~~~~~~~~~~~~~~~~~~~~~~~~~~~~~~~~~
\begin{equation*}
    \Gamma\ll1
    \qquad\Longleftrightarrow\qquad
    N_D\gg1.
\end{equation*}
% ~~~~~~~~~~~~~~~~~~~~~~~~~~~~~~~~~~~~~~~~~~~~~~~~~~~~~~~~

Числові множники тут важливі: $N_D=(6\sqrt\pi)^{-1}\Gamma^{-3/2}\approx0{,}09\,\Gamma^{-3/2}$, тож $\Gamma=0{,}07$ дає лише
$N_D\approx5$. Через малий коефіцієнт на практиці знак \enquote{$\gg$} слід читати як $\Gamma\lesssim0{,}01$, тобто
$N_D\gtrsim100$; ядро Сонця ($\Gamma\approx0{,}07$, $N_D\approx5$)~--- межа застосовності колективного опису.

Це два способи описати один і~той самий режим: $\Gamma$ говорить про силу взаємодії між
сусідніми частинками, а~$N_D$~--- про кількість частинок, що бере участь у~колективному
екрануванні.

Для слабкозв'язаної плазми, крім того, справедливе наближення ідеального газу (поправка Дебая~--- Хюккеля: потенціал оболонки в~точці самої частинки відрізняється від кулонівського на $-q/\lambda_D$, тож кореляційна енергія $W_{\text{кор}}\approx-e^2/2\lambda_D$ на частинку, а~відносно $|W_{\text{кор}}|/T\sim e^2/(\lambda_DT)\sim1/N_D\ll1$; вона дає лише малу поправку до тиску):

% ~~~~~~~~~~~~~~~~~~~~~~~~~~~~~~~~~~~~~~~~~~~~~~~~~~~~~~~~
\begin{equation}\label{eq:ideal-pressure}
    p\simeq\sum_s n_s T_s.
\end{equation}
% ~~~~~~~~~~~~~~~~~~~~~~~~~~~~~~~~~~~~~~~~~~~~~~~~~~~~~~~~

Нарешті, умова колективності сама по собі ще не гарантує квазінейтральності всієї системи:
для цього потрібна ще умова~\eqref{eq:size-condition}, $\lambda_D\ll L$. Тоді на макроскопічних
масштабах, значно більших за $\lambda_D$, зарядові неоднорідності швидко згладжуються і
$n_e\simeq Zn_i$.

Отже, маємо три різні умови з~різним фізичним змістом:
% ========================================================
\begin{enumerate}
	\item $N_D\gg1$~--- колективність екранування;
	\item $\lambda_D\ll L$~--- квазінейтральність на масштабі системи;
	\item $\omega\tau\gg1$~--- плазма поводиться як плазма, а~не як нейтральний газ ($\omega$~--- характерна частота процесу): за час $1/\omega$
	  заряджена частинка встигає відчути колективне поле, а~не лише зіткнутися з~нейтралом ($\tau$~--- середній
	  час між зіткненнями заряджених частинок з~нейтральними атомами). У~слабко іонізованому газі з~частими
	  зіткненнями з~нейтралами ця умова порушується, і~газ веде себе як звичайний, хоч $N_D\gg1$. Докладніше~---
	  в~розд.~\ref{sec:plasma-waves}.
\end{enumerate}
% ========================================================

Саме ці властивості дають змогу розглядати плазму як колективне квазінейтральне
середовище.

%% --------------------------------------------------------
\section{Нелінійне екранування біля зарядженої поверхні: шар Гуї~--- Чепмена}\label{sec:debye-layer}
%% --------------------------------------------------------


Лінійна теорія екранування вимагає $|q_s\varphi|\ll T_s$. Біля стінок камери,
зондів і~електродів ця умова може порушуватися. Тоді сама ідея дебаївського масштабу
залишається корисною, але експоненціальна формула для потенціалу вже не є точною. Цей
приклад важливий передусім як нагадування: \emphx{дебаївська довжина є масштабом
колективної перебудови, а~не універсальною формою потенціалу}. Для плоскої геометрії
нелінійну задачу можна розв'язати точно.

Нехай півпростір $x>0$ заповнено плазмою з~електронів і~однозарядних іонів (концентрація
вдалині $n$, $T_e=T_i=T$), а~на межі $x=0$ розташована площина з~поверхневою густиною
заряду $\sigma$. За Больцманом концентрації $n^{\pm}=n\exp(\mp e\varphi/T)$,
тому густина заряду $\rho=e(n^+-n^-)=-2ne\sinh(e\varphi/T)$, а~рівняння
Пуассона дає ($E=-d\varphi/dx$)
% ~~~~~~~~~~~~~~~~~~~~~~~~~~~~~~~~~~~~~~~~~~~~~~~~~~~~~~~~
\begin{equation*}
	\frac{dE}{dx}=4\pi\rho=-8\pi ne\sinh\frac{e\varphi}{T}.
\end{equation*}
% ~~~~~~~~~~~~~~~~~~~~~~~~~~~~~~~~~~~~~~~~~~~~~~~~~~~~~~~~
Помножимо обидві частини на $2E$: ліворуч вийде $2E\,E'=(E^2)'$, а~праворуч, з~огляду на $E\,dx=-d\varphi$,
\begin{equation*}
	d(E^2)=16\pi ne\sinh\frac{e\varphi}{T}\,d\varphi.
\end{equation*}
Це \enquote{інтеграл енергії}, той самий прийом, що й~для
рівняння $\ddot x=F(x)$ в~механіці (множимо на $2\dot x$ і~отримуємо $\tfrac12\dot x^2+U(x)=\text{const}$),
з~$x$ у~ролі часу, $\varphi$~--- координати, а~$E^2$~--- кінетичної енергії. Інтегруємо з~умовою $E\to0$ при $\varphi\to0$
($x\to\infty$) і~користуємося тотожністю $\cosh y-1=2\sinh^2(y/2)$:
% ~~~~~~~~~~~~~~~~~~~~~~~~~~~~~~~~~~~~~~~~~~~~~~~~~~~~~~~~
\begin{equation*}
	E^2=16\pi nT\left(\cosh\frac{e\varphi}{T}-1\right)=32\pi nT\sinh^2\frac{e\varphi}{2T}.
\end{equation*}
% ~~~~~~~~~~~~~~~~~~~~~~~~~~~~~~~~~~~~~~~~~~~~~~~~~~~~~~~~
Знак кореня такий, що $E$ і~$\varphi$ мають однаковий знак (бо $E=-d\varphi/dx$, а~$|\varphi|$ спадає від стінки вглиб плазми до нуля), отже
% ~~~~~~~~~~~~~~~~~~~~~~~~~~~~~~~~~~~~~~~~~~~~~~~~~~~~~~~~
\begin{equation*}
	-\frac{d\varphi}{dx}=4\sqrt{2\pi nT}\,\sinh\frac{e\varphi}{2T}.
\end{equation*}
% ~~~~~~~~~~~~~~~~~~~~~~~~~~~~~~~~~~~~~~~~~~~~~~~~~~~~~~~~
Це рівняння з~відокремлюваними змінними. З~$y=e\varphi/T$ воно має вигляд
$dy/\sinh(y/2)=-2\,dx/\lambda_D$, а~інтеграл
% ~~~~~~~~~~~~~~~~~~~~~~~~~~~~~~~~~~~~~~~~~~~~~~~~~~~~~~~~
\begin{equation*}
	\int\frac{dy}{\sinh(y/2)}=2\ln\tanh\frac y4
\end{equation*}
% ~~~~~~~~~~~~~~~~~~~~~~~~~~~~~~~~~~~~~~~~~~~~~~~~~~~~~~~~
легко перевірити диференціюванням ($\sinh(y/2)=2\sinh(y/4)\cosh(y/4)$). Його розв'язок
% ~~~~~~~~~~~~~~~~~~~~~~~~~~~~~~~~~~~~~~~~~~~~~~~~~~~~~~~~
\begin{equation}\label{eq:gouy-chapman}
	\tanh\frac{e\varphi}{4T}=\tanh\frac{e\varphi_0}{4T}\,e^{-x/\lambda_D},\qquad
	\lambda_D=\sqrt{\frac{T}{8\pi ne^2}},
\end{equation}
% ~~~~~~~~~~~~~~~~~~~~~~~~~~~~~~~~~~~~~~~~~~~~~~~~~~~~~~~~
де $\varphi_0=\varphi(0)$. Коефіцієнт в~експоненті збігається з~дебаївською
довжиною~\eqref{eq:debye-radius} (з~урахуванням обох сортів при $T_e=T_i=T$, тобто
$\lambda_D=\lambda_{De}/\sqrt2$), тобто вона є точним масштабом і~у~нелінійному режимі.
Зв'язок $\varphi_0$ з~$\sigma$ дає теорема Гаусса: усередині плазми екрануючий заряд
компенсує $\sigma$, поза площиною поля немає, тому $E(0)=4\pi\sigma$, звідки
% ~~~~~~~~~~~~~~~~~~~~~~~~~~~~~~~~~~~~~~~~~~~~~~~~~~~~~~~~
\begin{equation*}
	\sinh\frac{e\varphi_0}{2T}=\sigma\sqrt{\frac{\pi}{2nT}}.
\end{equation*}
% ~~~~~~~~~~~~~~~~~~~~~~~~~~~~~~~~~~~~~~~~~~~~~~~~~~~~~~~~
Два граничні випадки. При $e\varphi_0\ll T$ маємо
$\varphi=\varphi_0e^{-x/\lambda_D}$ з~$\varphi_0=4\pi\sigma\lambda_D$, як і~в~лінійній
теорії. При $e\varphi_0\gg T$ права частина~\eqref{eq:gouy-chapman} прямує до
одиниці, і~поблизу поверхні ($x\ll\lambda_D$) $\varphi\approx(2T/e)\ln(2\lambda_D/x)$: потенціал зростає до стінки лише логарифмічно, а~товщина шару
залишається порядку $\lambda_D$. Саме такі \emph{дебаївські шари} утворюються біля стінок
і~зондів, де квазінейтральність порушена. Порівняння точного розв'язку з~лінійною теорією
для $e\varphi_0=6T$ показано на рис.~\ref{pic:gouy-chapman}: далеко від стінки обидва спадають з~тим самим масштабом $\lambda_D$,
але нелінійний потенціал менший: його амплітуда $(4T/e)\tanh(e\varphi_0/4T)$ менша за $\varphi_0$ (потенціал \enquote{насичується}).

\begin{figure}[ht]\centering
	\localinput{gouy-chapman.tikz}
	\caption{Потенціал біля зарядженої площини при $e\varphi_0=6T$ (вісь $\varphi$ логарифмічна): точний розв'язок
	Гуї~--- Чепмена~\eqref{eq:gouy-chapman} і~лінійна теорія $\varphi_0e^{-x/\lambda_D}$. Далеко від стінки спад експоненціальний
	із~тим самим $\lambda_D$, але з~меншою амплітудою; біля стінки~--- логарифмічний}
	\label{pic:gouy-chapman}
\end{figure}


%% --------------------------------------------------------
\section{Шар біля стінки: критерій Бома і~зонд Ленгмюра}\label{sec:bohm-sheath}
%% --------------------------------------------------------

Розв'язок Гуї~--- Чепмена описує шар у~рівновазі з~Больцмановим розподілом \emph{обох}
сортів частинок. Біля реальної стінки, яка поглинає частинки, ситуація інша: електрони
значно швидші за іони ($v_{Te}/v_{Ti}\sim\sqrt{m_i/m_e}$), тож спочатку вони досягають
стінки й~заряджають її негативно відносно плазми. Утворюється шар додатного просторового
заряду, який відштовхує електрони й~прискорює іони (рис.~\ref{pic:bohm-sheath}). Цей
шар визначає потоки на стінку, а~отже, і~навантаження на першу стінку термоядерного
реактора, і~роботу зондів.

\begin{figure}[ht]\centering
	\localinput{sheath.tikz}
	\caption{Профіль потенціалу біля поглинальної стінки: квазінейтральний передшар, де іони
	набирають швидкість $\gtrsim c_s$, і~дебаївський шар, у~якому $n_i>n_e$}
	\label{pic:bohm-sheath}
\end{figure}

\subsection*{Критерій Бома}
\emph{Чому іони мають прийти в~шар прискореними.} Нехай $u$~--- швидкість іонів у~напрямку стінки.
При русі до стінки потенціал $\varphi$ спадає, іони прискорюються, а~їхня концентрація зменшується (збереження потоку
$n_iu=\mathrm{const}$). Електрони відштовхуються стінкою, і~їхня концентрація зменшується за Больцманом, $\propto\exp(e\varphi/T_e)$.
Щоб шар мав позитивний об'ємний заряд ($n_i>n_e$), густина іонів має спадати \emph{повільніше} за густину електронів.
Якщо іон входить у~шар повільно, то в~слабкому полі він різко прискорюється, і~його густина падає швидше, ніж електронна, а~отже,
виникає надлишок \emph{негативного} заряду, який не дає шарові існувати. Тому іони мають
прийти до шару вже з~помітною швидкістю (рис.~\ref{pic:bohm-profiles}). Кількісно цю умову дає такий розрахунок.

\begin{figure}[ht]\centering
	\localinput{bohm-profiles.tikz}
	\caption{Концентрації біля стінки (схематично): у~передшарі $n_i=n_e$, у~шарі електрони відштовхуються
	і~їхня концентрація $n_e$ падає за Больцманом швидше, ніж концентрація прискорених іонів $n_i$, тож $n_i>n_e$}
	\label{pic:bohm-profiles}
\end{figure}

Відлічуватимемо потенціал від межі шару ($\varphi=0$ там, де починається область $n_i\ne n_e$; концентрацію на
межі шару позначимо $n_\mathrm{sh}$). Нехай холодні іони входять туди зі~швидкістю $u_0$ в~напрямку стінки, а~електрони розподілені
за Больцманом. Із~закону збереження енергії $\tfrac12m_iu^2=\tfrac12m_iu_0^2-e\varphi$ та
потоку $n_iu=n_\mathrm{sh}u_0$ для іонів
$n_i=n_\mathrm{sh}\bigl(1-2e\varphi/m_iu_0^2\bigr)^{-1/2}$, а~$n_e=n_\mathrm{sh}\exp(e\varphi/T_e)$.
Введемо безрозмірний потенціал $\chi=-e\varphi/T_e>0$ і~параметр
% ~~~~~~~~~~~~~~~~~~~~~~~~~~~~~~~~~~~~~~~~~~~~~~~~~~~~~~~~
\begin{equation}\label{eq:ion-sound-speed-bohm}
	c_s=\sqrt{\frac{T_e}{m_i}}\quad(T_i\ll T_e,\ Z=1),
\end{equation}
% ~~~~~~~~~~~~~~~~~~~~~~~~~~~~~~~~~~~~~~~~~~~~~~~~~~~~~~~~
який має розмірність швидкості й~називається \emphx{іонно-звуковою швидкістю} (чому саме так, з'ясуємо в~розд.~\ref{sec:ion-acoustic}).
Тоді $n_i=n_\mathrm{sh}\bigl(1+2\chi c_s^2/u_0^2\bigr)^{-1/2}\approx n_\mathrm{sh}\bigl(1-\chi c_s^2/u_0^2\bigr)$ і~$n_e=n_\mathrm{sh}e^{-\chi}\approx n_\mathrm{sh}(1-\chi)$ при малих $\chi$, звідки
$n_i-n_e\approx n_\mathrm{sh}\chi\,(1-c_s^2/u_0^2)$.
Рівняння Пуассона $\varphi''=-4\pi e(n_i-n_e)$ у~змінній $\chi$ набуває вигляду
$\chi''=\chi\,(1-c_s^2/u_0^2)/\lambda_{De}^2$, де $\lambda_{De}^2=T_e/4\pi n_\mathrm{sh}e^2$.
Якби $1-c_s^2/u_0^2<0$, це було б рівняння гармонічного осцилятора: потенціал коливався б,
а~шар має монотонно відокремлювати плазму від стінки. Отже, на самій межі шару, де $\chi$ ще мале (тому лінеаризації досить), потрібен експоненціальний, а~не коливний розв'язок
(монотонне наростання потенціалу в~бік стінки); він можливий лише за умови
% ~~~~~~~~~~~~~~~~~~~~~~~~~~~~~~~~~~~~~~~~~~~~~~~~~~~~~~~~
\begin{equation}\label{eq:bohm-criterion}
	\tcbhighmath{u_0\ge c_s=\sqrt{T_e/m_i}.}
\end{equation}
% ~~~~~~~~~~~~~~~~~~~~~~~~~~~~~~~~~~~~~~~~~~~~~~~~~~~~~~~~
Це \emphx{критерій Бома}: іони мають увійти в~шар зі~швидкістю не меншою за
іонно-звукову. Прискорюються вони в~квазінейтральному \emph{передшарі}: там $n_i=n_e$ з~високою точністю, але потенціал
усе ж повільно спадає на масштабі $L\gg\lambda_D$ (довжина пробігу або розмір плазми). Електрони за Больцманом
\enquote{відчувають} це спадання, а~іони, що зникають на стінці, втягуються полем у~потік до неї.
Оцінимо перепад у~граничному випадку $u_0=c_s$: іони, що почали рух із~малої швидкості, набирають
$\tfrac12m_iu_0^2=\tfrac12T_e$, тож падіння потенціалу в~передшарі $\approx T_e/2e$, а~електрони за Больцманом
дають $n_\mathrm{sh}=n_0\exp(-1/2)\approx0{,}61\,n_0$ (квазінейтральність: $n_i=n_e$), де $n_0$~--- концентрація
в~незбуреній плазмі, далеко від стінки. Це оцінка, а~не строге виведення: реальний перепад залежить від того,
як іони народжуються й~гальмуються в~передшарі (іонізація, зіткнення), і~не менший за $T_e/2e$.

\subsection*{Плаваючий потенціал і~потік енергії на стінку}
Ізольований зонд (чи стінка) набуває потенціалу $\varphi_f$ відносно плазми, за якого
сумарний струм дорівнює нулю. Тепер потенціал відлічуємо від незбуреної плазми (на відміну від виводу критерію Бома,
де нуль стояв на межі шару), тож перепад $T_e/2e$ передшару входить у~$\varphi_f$. Потік іонів на стінку
$j_i=n_\mathrm{sh}c_s=n_0\exp(-1/2)\,c_s$. Електрони розподілені за Больцманом у~всьому перепаді від плазми
до стінки, тому до стінки доходить потік $j_e=\tfrac14n_0\bar v_e\exp(e\varphi_f/T_e)$, де
$\bar v_e=\sqrt{8T_e/\pi m_e}$ (однобічний потік $\tfrac14n\bar v$ нагадуємо з~молекулярної фізики: $n\int_{v_x>0}v_xf\,d^3v=n\bar v/4$). Прирівнюючи, маємо
% ~~~~~~~~~~~~~~~~~~~~~~~~~~~~~~~~~~~~~~~~~~~~~~~~~~~~~~~~
\begin{equation}\label{eq:floating-potential}
	\varphi_f=-\frac{T_e}{2e}\left(\ln\frac{m_i}{2\pi m_e}+1\right),
\end{equation}
% ~~~~~~~~~~~~~~~~~~~~~~~~~~~~~~~~~~~~~~~~~~~~~~~~~~~~~~~~
що для водню дає $\varphi_f\approx-3{,}3\,T_e/e$, для дейтерію~--- $\approx-3{,}7\,T_e/e$.
Із~цього перепаду $T_e/2e$ припадає на передшар, решта ($\approx2{,}8\,T_e/e$ для водню)~--- на
дебаївський шар. Іон, пройшовши весь перепад, приносить на стінку енергію $|e\varphi_f|$, тобто
кілька $T_e$. Разом з~електронами (кожен, що подолав бар'єр, приносить в~середньому $\approx2T_e$) і~початковою
тепловою енергією іонів ($\approx2{,}5T_i$) повний потік енергії на стінку $q\approx\gamma T_e\,j_i$ з~множником
$\gamma\approx2+e|\varphi_f|/T_e+2{,}5\,T_i/T_e$; для водню при $T_i\sim T_e$ це $\gamma\approx7$--$8$.

\subsection*{Товщина шару (для допитливих)}
Оцінимо товщину дебаївського шару $d$. Нехай у~шарі є лише іони (електрони відштовхнуто),
які для простоти стартують з~нуля (насправді з~$c_s$, але $\tfrac12m_ic_s^2=\tfrac12T_e\ll|e\varphi_s|$, тож це лише невелика поправка) і~падають у~потенціалі $\varphi(x)$, відлічуваному від межі шару ($x=0$) у~бік стінки і~взятому за модулем. Збереження енергії дає
$v=\sqrt{2e\varphi/m_i}$, а~стала густина струму $j=en_iv$ дає $n_i=j/(ev)$. Рівняння Пуассона
$\varphi''=4\pi e n_i=4\pi j\sqrt{m_i/2e}\;\varphi^{-1/2}$ має розв'язок $\varphi=Cx^{4/3}$
(підстановкою: $\tfrac49C^{3/2}=4\pi j\sqrt{m_i/2e}$), тобто
$\varphi\propto x^{4/3}$ (закон Чайлда~--- Ленгмюра). Якщо падіння потенціалу в~шарі $|\varphi_s|$ досягається при $x=d$, то
$C=|\varphi_s|/d^{4/3}$ і
\begin{equation*}
	j=\frac{1}{9\pi}\sqrt{\frac{2e}{m_i}}\,\frac{|\varphi_s|^{3/2}}{d^2},
	\qquad\text{звідки}\qquad d\propto|\varphi_s|^{3/4}\ \text{при заданому } j.
\end{equation*}
Підставивши $j=e\,n_\mathrm{sh}c_s$ і~$|\varphi_s|=\xi\,T_e/e$, отримуємо $d\approx0{,}8\,\xi^{3/4}\lambda_{De}$
(з~точністю до того, що $n_\mathrm{sh}\approx n$). Для водню $\xi\approx2{,}8$ (з~попереднього абзацу), тож $d\approx1{,}7\,\lambda_{De}$: товщина шару
зростає із~падінням потенціалу повільно, а~за порядком величини лишається $\lambda_{De}$.

\subsection*{Зонд Ленгмюра}
Нехай зонд~--- малий електрод площею $S$ і~характерним розміром $r_p$, потенціал якого $V$
відраховують від опорного електрода (наприклад, заземленої камери). Потенціал самої плазми
щодо опорного електрода позначимо $V_p$, плаваючий потенціал зонда~--- $V_f$; тоді
$\varphi_f=V_f-V_p<0$~--- та сама величина, що в~формулі~\eqref{eq:floating-potential}.
Змінюючи $V$, вимірюють струм $I$ на зонд; додатним вважаємо електронний струм (з~плазми на зонд),
а~іонний струм має від'ємний знак. Вольт-амперна
характеристика (рис.~\ref{pic:probe-iv}) має три ділянки.
\begin{figure}[ht]\centering
	\localinput{probe-iv.tikz}
	\caption{Вольт-амперна характеристика одиночного зонда Ленгмюра ($V_p$~--- потенціал плазми,
	$V_f$~--- плаваючий потенціал)}
	\label{pic:probe-iv}
\end{figure}
\begin{enumerate}
	\item $V\ll V_f$: струм іонного насичення
	  $I_{is}\approx0{,}6\,e\,n_0c_sS$ (коефіцієнт $0{,}6\approx\exp(-1/2)$ дає $n_\mathrm{sh}/n_0$; $S$~--- площа зонда);
	\item $V_f<V<V_p$: електрони гальмуються. До зонда долітають лише ті, чия енергія перевищує бар'єр $e(V_p-V)$;
		  їхня частка в~максвеллівському розподілі $\exp[-e(V_p-V)/T_e]$, тому $I_e\propto\exp\bigl(e(V-V_p)/T_e\bigr)$.
		  Звідси з~нахилу $\ln(I+I_{is})$ як функції $V$ знаходять $T_e$: $d\ln I_e/dV=e/T_e$;
	\item $V\ge V_p$: електронне насичення, $I_{es}=\tfrac14en_0\bar v_eS$.
\end{enumerate}
Отже, одиночний зонд дає $T_e$ (нахил), $V_p$ (злам) і~$n_0$ (з~$I_{is}$ при відомому $T_e$)~---
локальні параметри плазми, яких не знайти іншими простими методами.

\emph{Числовий ескіз} (водень, $T_e=3$~еВ, $n=10^{10}~\mathrm{см}^{-3}$): $c_s\approx1{,}7\times10^{6}$~см/с;
$I_{is}/S\approx0{,}6\,en c_s\approx1{,}6$~мА/см$^2$; $I_{es}/S=\tfrac14en\bar v_e\approx46$~мА/см$^2$; $\varphi_f\approx-3{,}3\,T_e/e\approx-10$~В.
Нахил $\ln I(V)$ дає $1/T_e$ (больцманівська частина електронного струму), а~$I_{is}$ не залежить від $V$, бо критерій Бома фіксує потік іонів.

\subsection*{Подвійний зонд і~межі застосовності}
Одиночний зонд потребує опорного електрода, що торкається плазми, а~розряд з~високочастотними
коливаннями потенціалу спотворює характеристику. Цього уникають \emphx{подвійним зондом}~---
двома однаковими електродами, між якими прикладено напругу $V$, а~вся система плаває.
Струм у~колі замкнений через плазму: що тече в~один електрод, те витікає з~іншого. Нехай
електроди мають потенціали $V_1$, $V_2$ відносно плазми, $V=V_1-V_2$. Кожен збирає
іонний струм насичення $I_{is}$ (однаковий, бо електроди однакові) і~електронний
$I_{ek}\propto\exp(eV_k/T_e)$ (як у~ділянці~2 характеристики одиночного зонда), тож
$I=I_{e1}-I_{is}=-(I_{e2}-I_{is})$. Звідси $I_{e1}+I_{e2}=2I_{is}$ і~$I_{e1}/I_{e2}=\exp(eV/T_e)$,
а~$I=\tfrac12(I_{e1}-I_{e2})$. Підставивши, отримаємо
% ~~~~~~~~~~~~~~~~~~~~~~~~~~~~~~~~~~~~~~~~~~~~~~~~~~~~~~~~
\begin{equation}\label{eq:double-probe}
	I=I_{is}\,\frac{\exp(eV/T_e)-1}{\exp(eV/T_e)+1}=I_{is}\tanh\frac{eV}{2T_e}\,,\qquad
	T_e=\frac{e\,I_{is}}{2\,(dI/dV)_{V=0}}\,.
\end{equation}
% ~~~~~~~~~~~~~~~~~~~~~~~~~~~~~~~~~~~~~~~~~~~~~~~~~~~~~~~~
Струм не перевищує $I_{is}$, тому збурення плазми мале. Температура визначається нахилом характеристики в~нулі
(рис.~\ref{pic:double-probe}).

\begin{figure}[ht]\centering
	\localinput{double-probe.tikz}
	\caption{Характеристика подвійного зонда~\eqref{eq:double-probe}: симетрична, насичується на $\pm I_{is}$;
	нахил у~нулі дає $T_e$}
	\label{pic:double-probe}
\end{figure}

Зонд працює лише тоді, коли він не руйнує плазму, яку вимірює:
\begin{itemize}
	\item розмір зонда має бути значно меншим за довжину вільного пробігу $\lambda_{\mathrm{fp}}$
	(відстань між зіткненнями, розд.~\ref{sec:plasma-collisions}), інакше струм визначається зіткненнями, а~не шаром;
	\item розмір зонда має бути значно більшим за товщину шару: $r_p\gg\lambda_D$; лише тоді струм
	$I_{is}\approx0{,}6\,e\,n_0c_sS$ збирається з~площі $S$, а~шар можна вважати тонким;
	\item у~магнітному полі потрібно $r_{Le}\gtrsim r_p$, де $r_{Le}$~--- ларморівський радіус
	електрона (радіус кола, яким він обертається навколо силової лінії; формула~\eqref{eq:larmor}
	у~розд.~\ref{sec:particle-motion}). Практично $r_{Le}\approx2{,}4\sqrt{T_e[\text{еВ}]}\,/B[\text{Гс}]$~см. Якщо $r_{Le}\ll r_p$, електрони
	майже не рухаються впоперек поля й~потрапляють на зонд лише вздовж силових ліній: збір струму
	залежить від орієнтації зонда щодо $\vect B$, тобто \emph{анізотропний}, і~формули цього розділу не працюють;
	\item у~гарячій плазмі ($T_e\gtrsim50$~еВ, великі потоки) зонд перегрівається, тому в~токамаках його використовують лише на периферії;
	\item за наявності швидких електронів (розподіл не максвеллівський) нахил $\ln I(V)$ зменшується на великих $V$, і~$T_e$ виходить завищеною.
\end{itemize}


%% --------------------------------------------------------
\section{Екранування в~виродженому газі: довжина Томаса~--- Фермі}\label{sec:thomas-fermi}
%% --------------------------------------------------------


Виведення дебаївського екранування спирається не на закон Больцмана як такий, а~лише на
чутливість концентрації електронів до потенціалу. Нехай електрони перебувають у~рівновазі
з~хімічним потенціалом $\mu$ (нерухомий іонний фон). У~потенціалі $\varphi$ електрон має
додаткову енергію $-e\varphi$, тому рівновага вимагає, щоб локальна величина
$\mu_{\mathrm{лок}}(n)-e\varphi$ була сталою. Звідси зміна концентрації $\delta
n=(\partial n/\partial\mu)\,e\varphi$, густина заряду $\rho=-e\,\delta n$, і~рівняння
Пуассона набуває вигляду $\Delta\varphi=\varphi/\lambda^2$ з

% ~~~~~~~~~~~~~~~~~~~~~~~~~~~~~~~~~~~~~~~~~~~~~~~~~~~~~~~~
\begin{equation}\label{eq:screening-general}
	\frac1{\lambda^2}=4\pi e^2\,\frac{\partial n}{\partial\mu}.
\end{equation}
% ~~~~~~~~~~~~~~~~~~~~~~~~~~~~~~~~~~~~~~~~~~~~~~~~~~~~~~~~

Для невиродженого газу $n\propto e^{\mu/T}$, $\partial
n/\partial\mu=n/T$, і~одержуємо $\lambda_{De}$. Для виродженого електронного
газу металу ($T\ll E_F$) $n\propto\mu^{3/2}$, тобто $\partial
n/\partial\mu=3n/2E_F$, і
% ~~~~~~~~~~~~~~~~~~~~~~~~~~~~~~~~~~~~~~~~~~~~~~~~~~~~~~~~
\begin{equation}\label{eq:thomas-fermi}
	\lambda_{TF}=\sqrt{\frac{E_F}{6\pi ne^2}}\approx0{,}5\sqrt{a_B\,n^{-1/3}},\qquad a_B=\frac{\hbar^2}{me^2}=0{,}53\times10^{-8}~\text{см},
\end{equation}
% ~~~~~~~~~~~~~~~~~~~~~~~~~~~~~~~~~~~~~~~~~~~~~~~~~~~~~~~~~~~~~~~
де використано $E_F=(\hbar^2/2m)(3\pi^2n)^{2/3}$. Тобто в~квантовій плазмі металу тепловий
масштаб $T$ замінюється енергією Фермі, а~екранування визначається середньою
відстанню між електронами. Для міді ($n\approx8{,}5\times10^{22}~\mathrm{см}^{-3}$,
$E_F\approx7$~еВ) $\lambda_{TF}\approx0{,}55\times10^{-8}$~см: заряд провідника
зосереджений у~шарі завтовшки в~частці міжатомної відстані. Це й~пояснює, чому
в~електростатиці так добре працює наближення \enquote{заряд розподілений по поверхні}. Тут важлива не стільки сама
оцінка $\lambda_{TF}$, скільки ідея: \emphx{масштаб екранування визначається властивостями розподілу частинок}, а~не
самим фактом наявності заряду. Для плазми це не екзотика: у~щільних астрофізичних об'єктах (білі карлики, ядра планет) електрони вироджені, і~замість $\lambda_{De}$ працює $\lambda_{TF}$.

%% --------------------------------------------------------
\begin{Summary}
\begin{itemize}
\item Плазма~--- квазінейтральний частково чи повністю іонізований газ із~колективною поведінкою заряджених частинок.
\item За формулою Саха газ помітно іонізується вже за $T\ll I$: вільний електрон має величезний фазовий об'єм, і~виграш ентропії компенсує больцманівський множник $e^{-I/T}$.
\item Дебаївська довжина $\lambda_D$~--- масштаб екранування заряду; на відстанях $L\gg\lambda_D$ плазма квазінейтральна.
\item Електрони й~іони можуть мати різні температури ($T_e\gg T_i$ у~розрядах), бо при зіткненні електрон віддає іону лише частку $\sim m_e/m_i$ енергії.
\item Колективність потребує великого числа частинок у~дебаївській сфері, $N_D\gg1$; це те саме, що слабка зв'язаність $\Gamma\ll1$ ($N_D\approx0{,}09\,\Gamma^{-3/2}$).
\item Плазма~--- це ще й~$\lambda_D\ll L$ (квазінейтральність на масштабі системи) та~$\omega\tau\gg1$ (немає домінування зіткнень з~нейтралами).
\item Біля поглинальної стінки іони входять у~шар об'ємного заряду зі швидкістю $u\ge c_s$ (критерій Бома); на цьому ґрунтується зонд Ленгмюра.
\item У~виродженому газі (метали) тепловий масштаб $T$ замінюється енергією Фермі: екранування задає довжина Томаса~--- Фермі.
\end{itemize}
\end{Summary}

\Needspace{9\baselineskip}
\section*{Задачі}
%% --------------------------------------------------------

\begin{problem}\label{prb:SunTemperature}
	Виразіть температуру поверхні Сонця ($T\approx5800$~К) в~електронвольтах, ергах і~джоулях.
	\emph{Відповідь:} $T\approx0{,}5$~еВ $\approx8\times10^{-13}$~ерг $\approx8\times10^{-20}$~Дж.
\end{problem}

\begin{problem}\label{prb:PlasmaProbe}
	Зонд Ленгмюра площею $S=0{,}1$~см$^2$ у~водневій плазмі з~$T_e=3$~еВ
	реєструє струм іонного насичення $1{,}6$~мА. Знайдіть концентрацію $n_0$ і~плаваючий
	потенціал відносно плазми.
	\emph{Відповідь:} $c_s\approx1{,}7\times10^{6}$~см/с, $n_0=I/(0{,}6\,e\,c_sS)\approx10^{11}~\mathrm{см}^{-3}$,
	$\varphi_f\approx-3{,}3\,T_e/e\approx-10$~В.
\end{problem}

\begin{problem}\label{prb:DoubleProbe}
	Подвійний зонд (формула~\eqref{eq:double-probe}) показує струм насичення $0{,}5$~мА,
	а~нахил характеристики в~нулі $dI/dV=0{,}1$~мА/В. Знайдіть $T_e$.
	\emph{Відповідь:} $T_e=eI_{is}/[2\,dI/dV]=0{,}5/(2\cdot0{,}1)=2{,}5$~еВ.
\end{problem}

\begin{problem}\label{prb:PlasmaSaha}
	За формулою Саха знайдіть ступінь іонізації водню з~повною концентрацією
	$10^{17}~\mathrm{см}^{-3}$ при $T=1$~еВ і~оцініть температуру, при якій
	$\alpha=1/2$.
	\emph{Відповідь:} $\alpha\approx0{,}18$; $\alpha=1/2$ при $T\approx1{,}2$~еВ.
\end{problem}

\begin{problem}\label{prb:PlasmaDebye}
	Знайдіть дебаївську довжину $\lambda_{De}$, число частинок у~дебаївській сфері $N_D$
	та параметр зв'язку $\Gamma$ для плазми з~$n_e=10^{12}~\mathrm{см}^{-3}$,
	$T_e=2$~еВ.
	\emph{Відповідь:} $\lambda_{De}\approx1{,}05\times10^{-3}$~см, $N_D\approx4{,}9\times10^{3}$, $\Gamma\approx7\times10^{-4}$.
\end{problem}

\begin{problem}\label{prb:PlasmaCollective}
	Для $n=10^{14}~\mathrm{см}^{-3}$, $T=10$~еВ знайдіть $\lambda_{De}$, $a$
	і~$N_D$. Перевірте співвідношення~\eqref{eq:ND-Gamma}. Чи виконується умова
	колективності?
	\emph{Відповідь:} $a\approx2{,}2\times10^{-5}$~см, $\lambda_{De}\approx2{,}4\times10^{-4}$~см ($\lambda_{De}/a\approx11$),
	$N_D\approx5\times10^{3}\gg1$; умова виконується.
\end{problem}

\begin{problem}\label{prb:PlasmaRelax}
	Оцініть час відновлення квазінейтральності $\tau\sim\lambda_{De}/v_{Te}$, $v_{Te}=\sqrt{T_e/m_e}$, для іоносфери
	та токамака з~табл.~\ref{tab:plasma-types}.
	\emph{Відповідь:} $\tau\approx6\times10^{-8}$~с (іоносфера), $\tau\approx2\times10^{-12}$~с (токамак).
\end{problem}

\begin{problem}\label{prb:PlasmaTwoTemp}
	Для плазми з~$n=10^{12}~\mathrm{см}^{-3}$, $T_e=3$~еВ, $T_i=0{,}03$~еВ знайдіть $\lambda_{De}$, $\lambda_{Di}$ і~$\lambda_D$.
	Яку довжину екранування слід брати для швидкого й~для повільного (порівняно з~$\lambda_D/v_{Ti}$) процесу?
	\emph{Відповідь:} $\lambda_{De}\approx1{,}3\times10^{-3}$~см, $\lambda_{Di}\approx1{,}3\times10^{-4}$~см, $\lambda_D\approx\lambda_{Di}$;
	для швидкого процесу~--- $\lambda_{De}$, для повільного~--- $\lambda_D$.
\end{problem}

\begin{problem}\label{prb:BohmFlux}
	Водневий розряд: $n_\mathrm{sh}=10^{11}~\mathrm{см}^{-3}$, $T_e=4$~еВ, $T_i\ll T_e$. Знайдіть мінімальну швидкість іонів на межі шару
	(критерій Бома), їхню енергію та потік на стінку.
	\emph{Відповідь:} $u_0=c_s\approx2{,}0\times10^{6}$~см/с, енергія $\tfrac12m_ic_s^2=T_e/2=2$~еВ,
	$j_i=n_\mathrm{sh}c_s\approx2\times10^{17}~\mathrm{см}^{-2}\mathrm{с}^{-1}$.
\end{problem}

\begin{problem}\label{prb:PlasmaDebyeLayer}
	Плазма з~$n=10^{12}~\mathrm{см}^{-3}$, $T=1$~еВ межує із~зарядженою
	площиною (розд.~\ref{sec:debye-layer}). Знайдіть $\lambda_D$ і~поверхневу густину заряду
	$\sigma$, при якій $e\varphi_0=T$. Покажіть, що при малих $\sigma$ маємо
	$\varphi_0=4\pi\sigma\lambda_D$.
	\emph{Відповідь:} $\lambda_D\approx5{,}3\times10^{-4}$~см, $\sigma\approx0{,}5$~СГС.
\end{problem}

\begin{problem}\label{prb:EdgePlasmaSheath}
	Пристінкова плазма токамака водню: $n=10^{13}~\mathrm{см}^{-3}$, $T_e=30$~еВ, $T_i\lesssim T_e$.
	(а)~Оцініть $\lambda_{De}$ і~$N_D$. (б)~Знайдіть плаваючий потенціал стінки відносно
	плазми за формулою~\eqref{eq:floating-potential}. (в)~Оцініть теплове навантаження на стінку
	$q\approx\gamma T_e\,j_i$, де $j_i=n_\mathrm{sh}c_s$, $n_\mathrm{sh}\approx\exp(-1/2)\,n\approx0{,}6\,n$
	(густина на межі шару), $\gamma\approx7$ (оцінка~--- у~тексті §\,\ref{sec:bohm-sheath}).
	\emph{Відповідь:} $\lambda_{De}\approx13$~мкм, $N_D\sim10^5$; $\varphi_f\approx-100$~В;
	$c_s\approx5{,}4\times10^6$~см/с, $j_i\approx3\times10^{23}~\mathrm{м}^{-2}\mathrm{с}^{-1}$,
	$q\sim10$~МВт/м$^2$.
\end{problem}

\begin{problem}\label{prb:PlasmaThomasFermi}
	Оцініть довжину Томаса~--- Фермі для міді ($n=8{,}5\times10^{22}~\mathrm{см}^{-3}$).
	Порівняйте з~міжатомною відстанню та поясніть, у~якому сенсі заряд провідника
	\enquote{розподілений по поверхні}.
	\emph{Відповідь:} $\lambda_{TF}\approx5\times10^{-9}$~см (поряд із міжатомною відстанню $\sim2\times10^{-8}$~см).
\end{problem}

%% --------------------------------------------------------
%% --------------------------------------------------------

\begin{Questions}
\begin{quizlist}
\item Чим плазма відрізняється від нейтрального газу? Що означає
		  \enquote{колективна} поведінка?

\item Чому газ іонізується помітно при температурі, багато меншій за енергію
		  іонізації? Що виражає формула Саха?

\item Чому в~газовому розряді $T_e\gg T_i$ і~що таке двотемпературна плазма?

\item Що таке квазінейтральність? Оцініть, наскільки вона порушується на масштабі
		  $L$.

\item Виведіть дебаївське екранування. Чому потрібен множник $e^{-r/\lambda_D}$
		  і~чому відкидається зростаючий розв'язок?

\item Чому в~слабкозв'язаній плазмі $N_D\gg1$? Як зв'язані $N_D$ і~$\Gamma$?

\item Чому дебаївський шар біля стінки має товщину $\sim\lambda_D$ навіть при
		  великому потенціалі?

\item Як змінюється довжина екранування у~виродженому електронному газі? Що таке
		  довжина Томаса~--- Фермі?

\item За яких умов плазма \enquote{колективна}? Назвіть три умови, їхній фізичний зміст і~співвідношення
		  між $N_D$ та $\Gamma$.

\item Чому іони мають увійти в~шар біля стінки зі~швидкістю $u\ge c_s$? Звідки береться плаваючий
		  потенціал $\approx-3\,T_e/e$ і~що вимірює зонд Ленгмюра?

\end{quizlist}
\end{Questions}
